\chapter{Related literature}\label{ch:lit}

This chapter places the thesis in four strands of work: the theory of why arbitrage is limited, the econometrics of the
band within which a futures basis moves, the evidence on short-sale constraints and prices, and the evidence from India.
It ends with what is new here.

\section{Why arbitrage is limited}
Shleifer and Vishny (1997) start from the fact that arbitrage is done by specialists who invest other people's money.
Investors judge them by past returns, so when a mispricing widens and the arbitrageur loses money, funds are withdrawn
exactly when the opportunity is best. In their three-period model the price can fall by more than one for one with a
deepening shock, because arbitrageurs are forced to sell into the fall. Two of their points matter here. Arbitrage
capital is limited and depends on the state of the market, and specialised arbitrageurs avoid positions whose returns are
volatile, so mispricing survives longest where idiosyncratic volatility is high.

De Long, Shleifer, Summers and Waldmann (1990) show that risk created by noise traders is enough to keep prices away from
fundamentals, even when the fundamental value is certain. Arbitrageurs with short horizons must sell at a price that
depends on future sentiment, and this risk deters them. The result needs the noise to be market-wide: misperceptions
that are specific to one asset can be diversified away by arbitrageurs who hold many such positions, so the model gives a
weaker reason to expect stock-level mispricing to rise with idiosyncratic volatility than Shleifer and Vishny do.
Loewenstein and Willard (2006) also show that the model's violation of the law of one price is not robust to changes in
its assumptions.

Gromb and Vayanos (2010) survey the theory and nest its main mechanisms in one model. A demand shock moves a price by an
amount that rises with the arbitrageurs' risk aversion and with the part of the asset's risk that cannot be hedged. A
cost of short selling adds a wedge only when arbitrageurs must be short, so the band around fair value becomes one-sided.
Margin requirements make capital a constraint, and the gap between two identical assets closes only when arbitrageurs
have enough of it. The model of Chapter~\ref{ch:theory} follows this structure and adds the feature that makes India
unusual: a position limit that, once reached, lets every participant only reduce positions.

Two classic papers explain why short-sale constraints matter for prices. Miller (1977) argues that when opinions differ
and pessimists cannot sell short, prices reflect the optimists and stocks are overpriced. Diamond and Verrecchia (1987)
show that with rational traders constraints do not bias prices on average but slow their adjustment to private
information, especially to bad news. Both predictions are tested in Chapter~\ref{ch:elig}.

\section{The no-arbitrage band in futures markets}
A futures contract and the underlying stock bought with borrowed money give the same payoff, so the futures price should
equal the cost of carry. Trading costs create a band around this value within which arbitrage does not pay. Brennan and
Schwartz (1990) add that an arbitrage position carries an option to close it before expiry if the gap reverses; the
model of this thesis uses the same idea.

Inside a band the basis should wander and outside it should revert, which calls for a threshold model. Dwyer, Locke and Yu (1996) estimate a three-regime
threshold error-correction model of index futures and cash. Using minute data on the S\&P 500 from 1982 to 1990, they find thresholds of about half a round-trip cost, faster
reversion outside the band than inside, and no asymmetry between the two sides. Tse (2001) argues that arbitrageurs differ
in their costs, so the aggregate response should rise smoothly with the deviation; his smooth-transition model of Dow
Jones futures fits well, and reversion survives when observations outside the band are removed, so it is not an artefact
of stale prices. Ters and Urban (2020) extend the threshold vector error-correction model of Hansen and Seo (2002) to three
regimes with an unknown cointegrating vector and a constant. They show that the threshold and the constant are each poorly
identified but their sum, which is the arbitrage cost, is estimated precisely, and they propose comparing the band before
and after regulatory changes. Chapter~\ref{ch:band} uses threshold and smooth-transition models in this tradition, on a
daily panel of stocks rather than on one index pair.

\section{Short-sale constraints and prices}
Ofek, Richardson and Whitelaw (2004) price stocks twice, once directly and once through put-call parity. Violations point
in the direction that short-sale costs predict: after shorting and trading costs, 13.6\% of their observations lie above
the upper bound and far fewer below it, violations are larger where borrowing is more expensive, and stocks that are expensive
relative to their options later earn low returns. Stambaugh, Yu and Yuan (2015) use the asymmetry of arbitrage to explain why
stocks with high idiosyncratic volatility earn low returns. With a mispricing score built from eleven anomalies, they show
in double-sorted portfolios that idiosyncratic volatility goes with lower returns among the most overpriced stocks
($-1.77\%$ a month for high minus low) and with higher returns among the most underpriced ($+0.52\%$), the first effect
being larger. The anomaly test of Chapter~\ref{ch:elig} follows their design with anomalies that can be built from public
Indian data.

Constraints can also be loosened. Gagnon (2018) studies about 1,300 introductions of single-stock futures in the United States and reports
that they lowered the use and the fees of securities lending, a sign of looser short-sale constraints. For the speed of price adjustment I use the price-delay measure of Hou and
Moskowitz (2005), which compares how much of a stock's response to market returns comes with a lag.

\section{Evidence from India}
The Indian basis has been studied mostly at the index level or in short intraday samples. Nandan, Agrawal and Bhargava
(2014) find that Nifty futures closed below their cost-of-carry value on 71\% of days from 2003 to 2012, with the gap
shrinking over the decade; they do not account for trading costs. Shankar, Sankar and Kiran Kumar (2019) estimate weekly
three-regime threshold models for 102 stock futures with five-minute data from 2007 to 2009. Both of their thresholds are
negative ($-39$ and $-10$\bp), the band is about 28\bp\ wide, it narrowed slightly when securities lending was introduced
in 2008, and the lower edge moves about twice as much as the upper with volatility, which points to a constraint on the
short side. Aggarwal (2015) estimates threshold models per stock and per day for 141 stock futures with one-second order
book data from March to August 2013. The average band runs from about $-15$ to $+3$\bp, it is wider and more asymmetric for
less liquid stocks, a common factor explains about 17\% of the variation in daily band widths, and that factor moved with
funding liquidity, for example when the Reserve Bank of India tightened liquidity in July 2013. Slivka, Shah and Wu (2011)
compute the full cost of arbitrage between a stock and its future and find a band about twice as wide as in mature
markets, with most arbitrage done within the day.

These bands are intraday objects for traders who close positions within minutes. The daily basis studied here measures the
part of mispricing that lasts overnight, where an arbitrageur must carry a position and pay delivery-based taxes, and so
is wider. Roy and Chakraborty (2023), using daily data for ten stocks from 2011 to 2017, find that stock and futures prices
are cointegrated and that the cash market does most of the correction of short-run deviations. Saurav, Srivastava and
Agarwalla (2026) study the stock-futures basis around earnings announcements. According to their abstract, the basis
turns negative before announcements and reverses afterwards, most strongly when disagreement is high and short-sale
constraints are tight; Chapter~\ref{ch:ban} finds the same pre-announcement pattern outside bans.

Two Indian studies of the cross-section bear on the anomaly tests. Harshita, Singh and Yadav (2018) find drift after
earnings announcements for Nifty 500 stocks from 2002 to 2017, measured with buy-and-hold returns from the second to the
sixty-fourth day after the announcement; the drift is not related to idiosyncratic volatility or illiquidity. Sehgal,
Vasishth and Deisting (2024) find that a stock's maximum daily return over the last five trading days of a month predicts low
returns the next month, but that high idiosyncratic volatility goes with \emph{higher} returns in India, the opposite of the United
States. Results from the United States therefore cannot be carried over to India without checking.

\section{Investors and regulation}
The demand that pushes stocks towards their position limits comes largely from speculation. Kumar (2009) shows that
individual investors in the United States prefer lottery-like stocks, which are cheap, volatile and positively skewed, and
that these stocks underperform. Barber and Odean (2000) show that the households that trade most earn the least: the most
active fifth earned 11.4\% a year net of costs against 17.9\% for the market. In India, SEBI's study of individual traders
in equity derivatives (SEBI, 2024a) found that 93\% of more than one crore individual traders lost money over the three
years to March 2024. SEBI changed the rules for index derivatives in October 2024, citing the rise in retail participation
and in speculative volumes (SEBI, 2024b), and changed the measurement of open interest and the position limits of single
stocks in May 2025 (SEBI, 2025); the second change is studied in Chapters~\ref{ch:band} and~\ref{ch:ban}.

\section{What this thesis adds}
The theory of limited arbitrage is mature, the band econometrics is standard, and Indian studies show that wide and
asymmetric bands exist. To my knowledge no study has estimated the effects of the F\&O ban, the rule under which every
participant may only reduce positions. The thesis models this rule, derives what it implies for the correction and the
level of the basis, and tests the predictions with a plan fixed in advance. It also estimates the band over a long daily
panel of stocks rather than an index or a short intraday sample, relates it to dated changes in costs, capital and
position limits, and uses entry into the F\&O segment and a constraint index for all listed stocks to study short-sale
constraints across the whole cash market.
